Estas bases no repiten el estado del arte. Toman de §\ref{sec:estado-arte} los papers que
\textbf{fundan} S, E o C, y los atan a OE1--OE4 y al Capítulo~4.

\subsection{Correspondencia estructural (\textit{schema matching})}

Shvaiko y Euzenat (2005) proponen una taxonomía multidimensional del \textbf{\textit{schema
matching}}: técnicas aproximadas o exactas a nivel de esquema; evidencia sintáctica, semántica o
externa a nivel de elemento y estructura; grado de automatización. El marco es anterior a OpenAPI y
JSON Schema, pero fundamenta el cálculo del \StructuralScore{} al contrastar \textit{paths},
\textit{operations} y \textit{schemas} con la organización de \Behavior{}s y \BusinessObject{}s del
\ServiceDomain{} de referencia.

\subsection{Alineación ontológica y similitud semántica}

El Bayadh et al.\ (2015) clasifican enfoques de \textbf{alineación ontológica} en cuatro familias:
medidas terminológicas, estructurales, extensionales y semánticas. Ninguna medida aislada basta;
combinaciones fijas pueden ser subóptimas, lo que justifica un diseño explícito de pesos en el
modelo S/E/C propuesto.

Chandrasekaran y Mago (2022) trazan la evolución de la \textbf{similitud semántica} desde recursos
léxicos (WordNet) hacia \textit{embeddings} y modelos transformadores. Santos Sousa et al.\ (2026)
catalogan el uso de \textit{embeddings} en ontology matching, señalando limitaciones de supervisión,
escalabilidad y correspondencias complejas.

Estas líneas aportan candidatos para el \SemanticScore{} y, en parte, para estimar cobertura
conceptual (\CoverageScore{}), siempre sobre vocabulario extraído del contrato OpenAPI y del modelo
BIAN normalizado.

\subsection{Representación enriquecida de contratos}

Mainas et al.\ (2023) ---citados en el eje E1--- proponen anotaciones semánticas
(\texttt{x-refersTo}, \texttt{x-kindOf}, etc.), ontología REST/Hydra y traducción OpenAPI
$\rightarrow$ RDF como vía de \textbf{representación enriquecida}.

Evaluaron más de 300 contratos de SwaggerHub con consultas SPARQL. Los tiempos de traducción superan
1~s para YAML de más de 200 líneas. En este plan se consideran como técnicas de normalización
candidatas (Cap.~4), no como definición previa del \SemanticScore{}.

\subsection{Relación con el modelo propuesto}

E2 y E3, solos, no son el producto. El grupo los mete en un procedimiento que entrega S, E, C y
\AlignmentScore{} sobre YAML bancario. Por qué hace falta esa unión: §\ref{sec:brecha}.
